\ifdefined\XeTeXrevision
\else
  \pdfoutput=1
\fi
\documentclass[11pt]{amsart}

\usepackage[T1]{fontenc}
\usepackage{lmodern}
\usepackage{microtype}
\usepackage{amsmath,amssymb,mathtools}
\usepackage{booktabs}
\usepackage[hidelinks]{hyperref}
\usepackage[nameinlink,capitalize,noabbrev]{cleveref}

\newtheorem{theorem}{Theorem}[section]
\newtheorem{proposition}[theorem]{Proposition}
\newtheorem{lemma}[theorem]{Lemma}
\newtheorem{corollary}[theorem]{Corollary}
\theoremstyle{remark}
\newtheorem{remark}[theorem]{Remark}

\DeclareMathOperator{\Spin}{Spin}
\DeclareMathOperator{\tr}{tr}
\DeclareMathOperator{\rank}{rank}
\DeclareMathOperator{\Gr}{Gr}

\title[Balanced Cayley information spectra]{Balanced Cayley Information Spectra for $\Spin(8)$ Triality Sensors}
\author{Hayden Austin}
\date{August 2026}
\keywords{$\Spin(8)$, triality, Cayley four-form, optimal design, information Gram operator}
\hypersetup{
  pdftitle={Balanced Cayley Information Spectra for Spin(8) Triality Sensors},
  pdfauthor={Hayden Austin},
  pdfsubject={Exact information spectra for orthonormal balanced Spin(8) triality sensors},
  pdfkeywords={Spin(8), triality, Cayley four-form, optimal design, information Gram operator}
}

\begin{document}

\begin{abstract}
Five infinitesimal observations of a shared $\Spin(8)$ action define a
$28\times28$ information Gram operator.  For the balanced allocation of one
vector, two positive-chiral, and two negative-chiral probes, the orthonormal
design space reduces to one Cayley parameter $c\in[-1,1]$.  We give an exact
block decomposition of the complete information family into fixed coordinate
subspaces of dimensions $8+8+8+4$ and prove
\[
  \det I_c=\frac{(1-c^2)^3(9-c^2)^2}{1024}.
\]
The direct spectral moments $\tr I_c=35$ and $\tr(I_c^2)=67$ are constant,
whereas $\tr(I_c^{-1})$ and $\tr(I_c^{-2})$ increase strictly with $c^2$.
Consequently the Cayley-null orbit is simultaneously D-optimal, A-optimal,
and inverse-Frobenius optimal within the complete orthonormal balanced family.
At the calibrated endpoints the rank is $25$ and the three lost eigenvalues
vanish at the common rate $(1-c^2)/8$.  The result is exact on the stated
orbit; it does not assert global optimality among nonorthogonal designs or
other allocations.
\end{abstract}

\maketitle

\section{Introduction}

Triality permutes the vector and two real chiral-spinor representations of
$\Spin(8)$, all of dimension eight.  This exceptional symmetry has long been
central to the geometry of the Cayley four-form and its $\Spin(7)$ stabilizer
\cite{HarveyLawson1982,KatzShnider2010}.  Here it is used as an information
geometry: a collection of probes observes one shared infinitesimal
$\mathfrak{spin}(8)$ action through the three triality-related views.

The resulting design problem is finite-dimensional but unusually rigid.  A
unit probe contributes a rank-seven orthogonal projector on the action
parameter space, which has dimension $28$.  Five probes are therefore the
first cardinality at which full local identifiability is possible.  The
balanced allocation studied here places one probe in the vector representation
and two in each chiral representation.

The main result is not only the determinant formula in the abstract.  A fixed
$8+8+8+4$ decomposition explains every determinant factor, proves an exact
duplication between two eight-dimensional blocks, and exposes a family whose
first two direct spectral moments remain fixed while its inverse moments
deteriorate monotonically.  This supplies three aligned design criteria and a
precise conditioning law at the singular endpoint.

The proof has two layers.  The standard cohomogeneity-one orbit description
for the $\Spin(7)$ action on oriented four-planes supplies the global orbit
parameter; it is recalled explicitly in the Cayley-hyperbolic-plane discussion
of \cite[Section~4]{BerndtTamaru2007}.  Exact rational calculations then verify
the internal $2+2$ flag quotient, construct the information operator, and
prove the block and spectral identities.  The computer-assisted portion is a
finite collection of exact symbolic identities, not a floating-point search.

\section{Triality information operators}

Let $V,S^+,S^-\cong\mathbb R^8$ denote the vector and two real chiral-spinor
representations.  Fix a bivector basis $E_1,\dots,E_{28}$ of
$\mathfrak{so}(8)$, and let $G_r^{(\alpha)}$ be the corresponding
skew-symmetric generator matrices in view $\alpha\in\{V,+,-\}$.

For a unit probe $x\in\mathbb R^8$, define
\[
 J_\alpha(x)
 =\begin{bmatrix}G_1^{(\alpha)}x&\cdots&G_{28}^{(\alpha)}x\end{bmatrix}
 \in\mathbb R^{8\times28},
 \qquad
 P_\alpha(x)=J_\alpha(x)^{\mathsf T}J_\alpha(x).
\]
The stabilizer of a unit vector is $\Spin(7)$, hence $P_\alpha(x)$ has rank
seven.  For a balanced design $(v;p_1,p_2;n_1,n_2)$, set
\[
 I=P_V(v)+P_+(p_1)+P_+(p_2)+P_-(n_1)+P_-(n_2).
\]

\begin{lemma}[same-view basis invariance]\label{lem:basis-invariance}
For an orthonormal pair $(x,y)$ in one triality view and
$\bigl(\begin{smallmatrix}r&t\\-t&r\end{smallmatrix}\bigr)\in O(2)$,
\[
 P_\alpha(rx+ty)+P_\alpha(-tx+ry)
 =P_\alpha(x)+P_\alpha(y).
\]
\end{lemma}

\begin{proof}
The Jacobian is linear in the probe.  Expanding both quadratic contributions
cancels the mixed terms and uses $r^2+t^2=1$.
\end{proof}

\section{The one-parameter balanced orbit}

The $\Spin(8)$ action first fixes $v=e_0$, leaving a $\Spin(7)$ stabilizer.
Clifford multiplication by $e_0$ identifies the two restricted chiral modules,
so the remaining probes determine an orthonormal four-frame in a common
$\mathbb R^8$.  The associated oriented four-plane is classified by its Cayley
coordinate $c\in[-1,1]$ under the cohomogeneity-one $\Spin(7)$ action; see the
explicit orbit description in \cite[Section~4]{BerndtTamaru2007} and the
triality/Cayley background in \cite{KatzShnider2010}.

The information operator retains a $2+2$ split of this four-plane.  The
following proposition records the additional flag statement required to show
that no second continuous invariant survives.

\begin{proposition}[balanced-flag normal form]\label{prop:normal-form}
Every orthonormal balanced design is information-spectrally equivalent to
\[
 (e_0;e_0,e_1;e_2,ce_3+se_4),
 \qquad c^2+s^2=1.
\]
After the independent $O(2)$ basis freedoms from
\cref{lem:basis-invariance}, the spectrum depends only on $z=c^2\in[0,1]$.
\end{proposition}

\begin{proof}
The classical orbit classification reduces the oriented four-plane to the
Cayley coordinate.  At an exact rational principal representative, the
four-plane stabilizer has dimension six and its restriction spans all of
$\mathfrak{so}(4)$.  Principal isotropy conjugacy propagates this statement
through the principal stratum.  Thus the stabilizer is transitive on oriented
orthogonal $2+2$ splittings, with stabilizer
$SO(2)\times SO(2)$.  Exact endpoint calculations give the same full
$\mathfrak{so}(4)$ restricted image on both singular orbits.  Finally,
reflecting one same-view pair changes $c$ to $-c$ while leaving its summed
information contribution unchanged.  The exact rational rank and restriction
identities are included in the accompanying certificate.
\end{proof}

\begin{remark}
The proposition deliberately separates the cited global four-plane orbit
classification from the exact internal flag calculation.  A local rank count
alone would not establish the global interval normal form.
\end{remark}

\section{Exact invariant blocks}

Let $I(c,s)$ denote the information operator of the representative in
\cref{prop:normal-form}.  All polynomial identities below are taken in
\(\mathbb Q[c,s]/(c^2+s^2-1)\).

\begin{theorem}[block characteristic laws]\label{thm:blocks}
A fixed permutation of the bivector coordinates gives
\[
 I(c,s)=I_8^{(0)}\oplus I_8^{(1)}\oplus I_8^{(2)}\oplus I_4.
\]
Writing $\lambda$ for the spectral variable, the monic block polynomials are
\begin{align*}
\chi_0(\lambda)
={}&-\frac14(\lambda-1)^2
 (2c\lambda-c-2\lambda^3+8\lambda^2-6\lambda+1)\\
&\hspace{37mm}\cdot
 (2c\lambda-c+2\lambda^3-8\lambda^2+6\lambda-1),\\
\chi_1(\lambda)=\chi_2(\lambda)
={}&\frac1{16}(c-2\lambda^2+4\lambda-1)
 (c-2\lambda^2+6\lambda-3)\\
&\hspace{20mm}\cdot(c+2\lambda^2-6\lambda+3)
 (c+2\lambda^2-4\lambda+1),\\
\chi_3(\lambda)={}&(\lambda-1)^2(\lambda^2-3\lambda+1).
\end{align*}
Moreover, $I_8^{(1)}$ and $I_8^{(2)}$ are conjugate by a fixed signed
permutation matrix.
\end{theorem}

\begin{proof}
The support graph of the exact $28\times28$ information matrix has four fixed
connected components.  Direct determinant calculation in the circle quotient
gives the displayed characteristic polynomials.  Substitution into the
expanded matrices leaves zero remainder.  A stored signed permutation $U$
satisfies $UI_8^{(1)}=I_8^{(2)}U$ and $U^{\mathsf T}U=I_8$ entrywise.
\end{proof}

\section{Determinant and aligned design criteria}

\begin{theorem}[balanced Cayley spectrum]\label{thm:main}
For $z=c^2\in[0,1]$,
\[
 \det I_c=\frac{(1-z)^3(9-z)^2}{1024}.
\]
The Cayley-null orbit $z=0$ is the unique D-optimum on the orthonormal balanced
family, with $\det I_0=81/1024$.  The endpoint $z=1$ has rank $25$.
\end{theorem}

\begin{proof}
Evaluating the block polynomials of \cref{thm:blocks} at zero gives
\[
 \det I_8^{(0)}=\frac{1-z}{4},\qquad
 \det I_8^{(1)}=\det I_8^{(2)}=\frac{(1-z)(9-z)}{16},
 \qquad \det I_4=1.
\]
Their product is the stated determinant.  Its derivative is
\[
 -\frac{(1-z)^2(9-z)(29-5z)}{1024},
\]
which is strictly negative for $0\le z<1$.  The exact endpoint block ranks
sum to $25$.
\end{proof}

\begin{theorem}[direct and inverse spectral moments]\label{thm:moments}
For every $z\in[0,1]$,
\[
 \tr I_c=35,
 \qquad
 \tr(I_c^2)=67.
\]
For $0\le z<1$,
\begin{align*}
 \tr(I_c^{-1})
 &=\frac{11z^2-206z+387}{(1-z)(9-z)},\\
 \tr(I_c^{-2})
 &=\frac{19z^4-76z^3+786z^2+2676z+8883}
 {(1-z)^2(9-z)^2}.
\end{align*}
Both inverse moments increase strictly with $z$.
\end{theorem}

\begin{proof}
The direct moments follow from the first coefficients of the block
characteristic polynomials.  For the inverse moments, logarithmic
differentiation of $p(\lambda)=\det(\lambda I-I_c)$ at zero gives the displayed
rational functions.  Their derivatives reduce respectively to
\[
 \frac{96(z^2-6z+21)}{(9-z)^2(1-z)^2}
\]
and
\[
 \frac{16(12609+336z-630z^2-8z^3-19z^4)}
 {(9-z)^3(1-z)^3}.
\]
The first numerator equals $96((z-3)^2+12)$.  On $[0,1]$, the second
polynomial is at least $12609-630-8-19>0$.
\end{proof}

\begin{corollary}
On this family, the Cayley-null orbit uniquely maximizes $\det I$ and uniquely
minimizes $\tr(I^{-1})$ and
$\tr(I^{-2})=\lVert I^{-1}\rVert_F^2$.
\end{corollary}

\section{Equal-rate loss of identifiability}

\begin{theorem}[endpoint asymptotics]\label{thm:endpoint}
As $z\uparrow1$, exactly three eigenvalues vanish and
\[
 \lambda_j(z)=\frac{1-z}{8}+O((1-z)^2),\qquad j=1,2,3.
\]
\end{theorem}

\begin{proof}
At $c=1$, each eight-dimensional block polynomial has a simple zero root.
Exact implicit differentiation gives $d\lambda_j/dc=-1/4$ for all three
blocks.  Since $1-c^2=-2(c-1)+O((c-1)^2)$, the common slope in $1-z$ is
$1/8$.
\end{proof}

\begin{corollary}[conditioning law]\label{cor:conditioning}
The largest surviving endpoint eigenvalue is $2+\sqrt2$, and
\begin{align*}
 \kappa_2(I_c)&=\frac{8(2+\sqrt2)}{1-z}+O(1),\\
 \tr(I_c^{-1})&=\frac{24}{1-z}+O(1),\\
 \tr(I_c^{-2})&=\frac{192}{(1-z)^2}+O((1-z)^{-1}).
\end{align*}
\end{corollary}

\begin{proof}
Factor the surviving endpoint polynomials in \cref{thm:blocks}, then invert
the three expansions from \cref{thm:endpoint}.
\end{proof}

\section{Proof objects and scope}

The accompanying repository constructs the three rational triality generator
families, checks their Lie brackets and common equivariance tensor, computes
the principal and endpoint flag stabilizers, builds the exact information
family, and verifies every block identity in the circle quotient.  Independent
SymPy and FLINT calculations replay selected derivative and Bernstein
identities \cite{Farouki2012}.  Artifact hashes establish byte provenance;
they do not replace the algebraic replay.

The theorem is confined to the complete orthonormal balanced orbit.  It does
not prove that orthonormalization improves every correlated frame, does not
compare all five-probe allocations, and does not establish any advantage for
sequence models.  Exact signed-star and residual-edge results address separate
nonorthogonal subfamilies, while the unrestricted Dirac--Gram inequality
remains open.

\section*{Data and code availability}

The exact source, tests, and machine-readable proof artifacts are included in
the associated repository.  The source archive records the Python and
python-flint versions, bounded thread settings, and SHA-256 manifest required
for replay.

\bibliographystyle{amsplain}
\bibliography{references}

\end{document}
